\subsubsection{Digit DP template}

\paragraph{Idea intuitiva.}
Contar numeros en $[0, N]$ que cumplen una propiedad sobre sus digitos. La idea: DP por posicion, manteniendo:
\begin{itemize}\setlength\itemsep{0pt}
	\item \texttt{pos}: posicion actual (de izquierda a derecha).
	\item \texttt{tight}: si los digitos ya colocados son exactamente los de $N$ (limite superior estricto).
	\item \texttt{state}: estado adicional (digitos usados, suma, etc.).
\end{itemize}
La clave: la condicion \texttt{tight} permite acotar correctamente el siguiente digito; sin ella, la DP contaria incorrectamente.

\paragraph{Propiedades clave (invariantes).}
\begin{itemize}\setlength\itemsep{0pt}
	\item Si \texttt{tight} es true, el siguiente digito esta acotado por el de $N$; si no, puede ser 0-9.
	\item Memoria: $O(D \cdot 2 \cdot \text{states})$.
	\item Caso base: si \texttt{pos == D}, devolver 1 si el estado cumple la propiedad, sino 0.
	\item El bit \texttt{started} distingue ``numero con menos digitos'' (con leading zeros).
\end{itemize}

\paragraph{Codigo clave.}
La estructura basica:
\begin{verbatim}
long long dp(int pos, bool tight, State s) {
    if (pos == D) return check(s) ? 1 : 0;
    if (memo[pos][tight][s] != -1) return memo;
    int limit = tight ? digits[pos] : 9;
    long long ans = 0;
    for (int d = 0; d <= limit; ++d)
        ans += dp(pos + 1, tight && (d == limit), update(s, d, pos));
    return memo[pos][tight][s] = ans;
}
\end{verbatim}

\paragraph{Complejidad.}
\begin{itemize}\setlength\itemsep{0pt}
	\item $O(D \cdot 10 \cdot \text{states})$ donde $D$ es el numero de digitos.
	\item Para $D \le 19$ (hasta $10^{18}$): $\sim 200 \cdot \text{states}$.
\end{itemize}

\paragraph{Donde tocar para modificar.}
\begin{itemize}\setlength\itemsep{0pt}
	\item \textbf{Cambiar la propiedad}: modificar \texttt{updateState} y la condicion base.
	\item \textbf{Multiples estados}: aniadir mas dimensiones al \texttt{memo}.
	\item \textbf{Leading zeros}: aniadir un flag \texttt{started} para distinguir numeros con menos digitos.
	\item \textbf{Rango $[L, R]$}: calcular $f(R) - f(L-1)$.
\end{itemize}

\paragraph{Trampas y errores frecuentes.}
\begin{itemize}\setlength\itemsep{0pt}
	\item Olvidar resetear \texttt{memo} entre queries.
	\item Si la cantidad de estados es grande, el \texttt{memo[pos][tight][state]} puede ser muy grande; usar \texttt{map}.
	\item La condicion \texttt{started} es crucial para numeros con menos digitos que $D$.
	\item Si $N$ es muy grande ($> 10^{18}$), aniadir soporte para long long.
\end{itemize}

\paragraph{Codigo.}
Ver \texttt{cpp.tex}, seccion \textit{Digit DP template}.

\vspace{0.5em}
